\section{Discussion}
\label{sec:discussion}

\paragraph{Why doesn't relevance convert?} Our decomposition suggests a
simple economy: a retrieved example helps only if it is \emph{sufficiently}
relevant to the specific fix, and realistic corpora rarely contain such an
example for a given new bug. When one exists (planted), even plain dense
retrieval finds it and repair converts; when none exists, improving the
\emph{ordering} of insufficient examples cannot manufacture the missing
knowledge. Reranking sophistication addresses the part of the problem that
was not binding.

\paragraph{Headroom scales with what the model currently fails.} The planted
gain is $+14$/$+13\pp$ on the two models with mid-range base accuracy and
$\sim 0$ on the saturated one. We frame this as ``headroom scales with the
fraction of currently-failed problems that a perfect example would fix''
rather than a categorical knowledge-gap principle: Qwen2.5-7B ($78.6\%$ base)
is close to GPT-4o's HumanEvalFix base ($80.5\%$), and the distinguishing
variable in our data is base accuracy, not model family.

\paragraph{Cost: structured RAG is Pareto-dominated for repair.} Measured on
HumanEvalFix GPT-4o traces, the structured prompt averages 4{,}632 input
characters ($\approx$1{,}158 tokens) against 1{,}199 ($\approx$300) for the
baseline --- a ${\sim}4\times$ input multiplier, plus embedding, indexing,
and retrieval infrastructure --- for equal-or-worse accuracy. In absolute
terms the marginal cost is small ($\sim$\$0.002/repair at GPT-4o input
pricing); the point is direction: you pay more and get nothing. The
actionable guidance for practitioners follows from RQ3: \textbf{spend the
retrieval budget on corpus coverage and curation, not reranking
sophistication} --- the two are separable, and only one showed positive
returns anywhere in this study.

\paragraph{Implications for reported RAG-repair gains.} Contamination and
post-processing can masquerade as retrieval effects. In our own PyBugHive
setting, an uncontrolled corpus (containing same-project entries) flipped
results from neutral to strongly negative-looking for RAG variants; the
project-held-out control removed the artifact. Externally, our recomputation
of RAGFix's released outputs (Section~\ref{sec:related}) is consistent with a
post-processing artifact. We recommend that RAG-repair evaluations
(i)~hold out same-project corpus entries, (ii)~ablate post-processing
separately from retrieval, and (iii)~report equivalence bounds, not bare
non-significance.
